\section{Hopf--Cole linearization and path-integral control}\label{sec:mppi}

\begin{remark}[Rigid-body configuration spaces and spectral truncation]
The configuration space of a floating-base rigid body with $n$ revolute joints is $Q = SE(3) \times \bT^n$. The Peter--Weyl basis for $L^2(G_1 \times G_2)$ consists of tensor products $\rho_1 \otimes \rho_2$ with Casimir eigenvalue $\lambda_{\rho_1 \otimes \rho_2} = \lambda_{\rho_1} + \lambda_{\rho_2}$.

Two obstructions arise in high dimension. First, $SE(3)$ is non-compact: its Plancherel decomposition has continuous spectrum and no spectral gap (\S\ref{sec:noncompact}). Second, on the compact factor $SO(3) \times \bT^n$, the Weyl eigenvalue asymptotics give $O(\Lambda^{(3+n)/2})$ spectral degrees of freedom below a cutoff~$\Lambda$---a curse of dimensionality for any spectral Galerkin truncation. Stochastic methods bypass this entirely: path sampling yields Monte Carlo estimates at rate $O(N^{-1/2})$, independently of $\dim G$---the CLT rate (\cref{thm:add-clt}). The spectral gap $\lambda_1$ (\cref{cor:spectral-gap}) re-enters as the mixing rate of the sampler.
\end{remark}

\begin{remark}[Hopf--Cole linearization and path-integral control]
Specialize the HJB equation (\S\ref{sec:hjb}) to quadratic control cost
$L(g,u) = q(g) + \tfrac{1}{2}\|u\|_{\mathfrak{g}}^2$.
The $\min_u$ is a Legendre--Fenchel transform: the minimizer is
$u^* = -\nabla_{\mathfrak{g}} V$ with minimum value
$q(g) - \tfrac{1}{2}\|\nabla_{\mathfrak{g}} V\|^2$.
The specialized HJB becomes
\[
\partial_t V - \varepsilon\,\Omega V + q(g)
  - \tfrac{1}{2}\|\nabla_{\mathfrak{g}} V\|^2 = 0,
\quad V(T,g) = \Phi(g).
\]
In the deterministic limit $\varepsilon = 0$, the characteristics are Euler--Poincar\'e trajectories on $\mathfrak{g}$ (equivalently, Lie--Poisson flow on $\mathfrak{g}^*$).

\medskip\noindent\textbf{Linearization.}
The Hopf--Cole substitution $\psi = e^{-V/(2\varepsilon)}$ exactly cancels the
quadratic gradient nonlinearity, yielding the linearized PDE:
\[
\partial_t \psi + (-\varepsilon\,\Omega)\psi
  - \frac{q}{2\varepsilon}\,\psi = 0,
\quad \psi(T,g) = e^{-\Phi(g)/(2\varepsilon)}.
\]
This is a backward linear PDE---the heat equation on $G$ with killing
potential $q/(2\varepsilon)$.

\medskip\noindent\textbf{Feynman--Kac representation.}
Since $-\varepsilon\,\Omega$ is the generator of the uncontrolled diffusion
$dg_s = g_s \cdot \sqrt{2\varepsilon}\,dW_s$:
\[
\psi(t,g) = \mathbb{E}_g\!\left[
  \exp\!\left(-\frac{1}{2\varepsilon}\int_t^T q(g_s)\,ds\right)
  e^{-\Phi(g_T)/(2\varepsilon)}
\right],
\]
where the expectation is over the uncontrolled (pure-diffusion) process
started at $g_t = g$.

\medskip\noindent\textbf{Optimal control recovery and MPPI.}
The optimal control is
$u^*(t,g) = 2\varepsilon\,\nabla_{\mathfrak{g}}\log\psi(t,g)$.
The nonlinear HJB is replaced by forward sampling of the SDE plus
exponential reweighting---this is model-predictive path-integral control
(MPPI) \cite{Kappen,Todorov}.
The Monte Carlo estimate converges at $O(N^{-1/2})$ for bounded costs, independently of $\dim G$---the CLT rate (\cref{thm:add-clt}).

\medskip\noindent\textbf{Deterministic limit.}
As $\varepsilon \to 0$, a Schilder-type large-deviation argument shows
$V(t,g) \to \min_u\bigl\{\int_t^T L(g_s,u_s)\,ds + \Phi(g_T)\bigr\}$.
The minimizing path satisfies the Euler--Lagrange equation; on a Lie group
with left-invariant kinetic energy, this reduces to the Euler--Poincar\'e
equation on $\mathfrak{g}$.
\end{remark}

\begin{remark}[Phase propagation and the eikonal limit]\label{rem:eikonal}
The deterministic limit $\varepsilon \to 0$ of the specialized HJB (preceding remark) yields the Hamilton--Jacobi equation $\partial_t V + H(\nabla V) = 0$. For wave propagation in a medium with velocity profile $v(x)$, the phase $\phi(x,t)$ satisfies the \emph{eikonal equation}:
\[
|\nabla\phi|^2 = \frac{1}{v(x)^2},
\]
which is precisely a Hamilton--Jacobi equation. The wavefronts are characteristics; the rays are geodesics on the Riemannian manifold $(M, g_{ij} = v(x)^{-2}\delta_{ij})$. This is exact---not an approximation, not sampling, not machine learning. Given the velocity structure, one computes arrival times at any point in one step.

\emph{Physical realizations.}
\begin{itemize}
\item \emph{Seismic waves.} Phase propagation through Earth's interior with velocity structure $v(x)$. Given the velocity profile, arrival times at any point are computed exactly via the eikonal equation.
\item \emph{Tsunamis.} Shallow water wave propagation: $\partial_{tt}\eta = \nabla\cdot(gH(x)\nabla\eta)$, with phase velocity $\sqrt{gH(x)}$ determined by known bathymetry $H(x)$. This is why tsunami early warning systems produce exact arrival-time predictions within minutes of an earthquake.
\item \emph{Elastic waves in kinematic chains.} Contact force propagation through a robot's structure---the same physics at smaller scale, with exactly known link geometry.
\end{itemize}

\emph{Peter--Weyl on $S^2$.}
Earth is (approximately) a sphere. The Peter--Weyl decomposition of $L^2(S^2)$ yields the spherical harmonics $Y_\ell^m$, and Earth's free oscillations (normal modes) are exactly these harmonics with eigenfrequencies $\omega_{n\ell}$ determined by the radial velocity structure. Measuring eigenfrequencies after large earthquakes and inverting for $v(x)$ is the inverse spectral problem---hearing the shape of the drum. The forward problem (given structure, compute modes) is exact because it reduces to solving Sturm--Liouville equations on each representation block, exactly as the heat equation block-diagonalizes via the Casimir spectrum (\cref{cor:spectral-gap}).

\emph{Cycle versus re-cycle.}
Event nucleation (when an earthquake occurs) may have infinite holonomy---chaotic, sensitive to initial conditions, possibly unpredictable. But once the source exists, propagation on known structure is exact computation, not sampling. The structure (velocity profile, bathymetry, link geometry) is the cycle; estimating it from observations (seismic tomography, sonar mapping, system identification) is the re-cycle.
\end{remark}

